\section{My team}

% ============================================================
% 2a. TEAM STRUCTURE
% ============================================================

\subsection{2a. Team structure --- discussed explicitly, and what changed}

At the wk12 kickoff we did hold an explicit structure conversation, and I want to credit Robin for forcing it: I would have skipped it. We split on module ownership --- Robin took the flight loop and state execution, I took perception and simulation, Edward took the Cube and hardware interface, Demetro took client-facing communication and the FDR presentation, and [PLACEHOLDER: fifth member name] coordinated schedule and project management. We deliberately declined to name a single technical lead because the brief (R08) devolves autonomy-level choices to the team and we wanted that authority distributed. We also agreed on a Git-based workflow: one main branch, feature branches merged only after informal peer review, and a shared \texttt{requirements\_pi.txt} that all five of us would update.

Mapping this against Tuckman's (1965) model, our Forming stage was brief --- arguably too brief. Graham warns in his Professional Practice lecture that ``too little time spent forming, which means no building of trust in the team, no discussion of its aims or its goals, and just jumping straight into the storming phase'' produces the worst-case scenario where the team eventually needs to reform under pressure. That is, I now recognise, approximately what happened to us.

The structure then drifted twice. The first drift was silent: around wk16, as hardware kept slipping and three flight days collapsed, the project's centre of gravity migrated from the bench to my laptop because that was the only surface where forward progress was possible. Nobody declared this shift; it just happened. I was running SITL simulation, writing test scripts, and building the web ground station while Edward waited for parts and Robin waited for flight data. The practical effect was that my role expanded from ``perception specialist'' to ``de facto technical infrastructure owner'' without anyone consciously choosing that. In Lencioni's (2002) terms, this is a textbook expression of Dysfunction 4, avoidance of accountability: ``without commitment to a clear plan, there is nothing to measure against --- low standards prevail.'' We had committed to a structure in wk12, but we had not committed to a mechanism for detecting when that structure no longer matched reality.

The second drift was loud: at wk17, after I caught Robin and myself building overlapping state machines in parallel --- I was adding states to \texttt{main.py} while Robin was adding equivalent logic to his own flight script --- I asked for a renegotiation. This was Tuckman's Storming stage, unmistakable: ``ideas start to come from individuals, rivals may emerge and compete for different work packages or roles'' (Graham, paraphrasing Tuckman). We agreed in writing (\texttt{docs/DESIGN\_DECISIONS.md:107--109}) that Robin's script would own the flight loop and call my \texttt{passive\_watch} as an HTTP service for detections. ``Robin's script handles flight only (no camera). Our \texttt{pi\_flight.py} handles camera + vision + stream. Both run in parallel, no conflict.'' That written contract was the boundary that stopped the overlap.

Looking back through Lencioni's pyramid, I can trace our dysfunction clearly. I never fully trusted anyone with my code (Level 1: absence of trust), which meant Robin's FSM pushback felt threatening rather than healthy (Level 2: fear of conflict), which meant we never achieved genuine commitment to a shared architecture until the wk17 crisis forced it (Level 3: lack of commitment). Lencioni's chain metaphor --- ``break one link and the whole chain fails'' --- describes our first sixteen weeks with uncomfortable accuracy.

One reading of this is that our team structure failed. Another is that it succeeded in the way that matters: we caught the overlap in wk17, not wk24, and the written contract held for the rest of the project. The evidence: Robin's wk20 integration of \texttt{passive\_watch} took an afternoon, not a week, and he did it without asking me a single clarifying question for two days. I conclude the structure worked not because it was well-designed at the start, but because we had the capacity to renegotiate when it broke. That capacity, I now believe, is more valuable than getting the initial design right.

\ebi\ We should have revisited the structure formally at least every two weeks, not only when it visibly broke. A standing agenda item --- ``does our current structure still match the work?'' --- would have caught the wk16 silent drift before it became the wk17 conflict.

% ============================================================
% 2b. MOST IMPACTFUL TEAM-WORKING ELEMENTS
% ============================================================

\subsection{2b. Most impactful elements of team working}

The most impactful team-working element, measured by what it actually unblocked, was not a tool or a practice. It was a \textit{handoff ratio} I had not previously noticed.

The four hours I spent writing \texttt{docs/VISION\_PIPELINE\_FOR\_COLLEAGUE.md} and the \texttt{--robin} integration flag (commit \texttt{a5b1ab7}) bought Robin roughly twenty hours of unblocked integration time across wk20. He wired \texttt{passive\_watch} into his mission script without asking me a single clarifying question for two full days. I now understand those docs worked because they traded my solo time for Robin's onboarding time at roughly 5:1 --- better team economics than I had assumed documentation ever delivered. I had previously treated handoff docs as a cost I paid for politeness. In fact they were the highest-leverage work I did on this project, measured by teammate-hours-unblocked per author-hour-invested.

\www\ The second most impactful element was the progressive test ladder. I organised all flight tests into numbered steps (\texttt{tests/flight/0a\_cube\_commands.py} through \texttt{4\_detect\_and\_center.py}), with a strict rule: never skip a step. Each step builds trust before adding risk. This was, in Graham's risk management terms, a mitigation that reduced both likelihood and consequence --- like the climbing rope, not the helmet. The test ladder meant that when we finally ran on the Pi on 2026-03-31, we knew which components had been verified and which hadn't. The delta between simulation and hardware turned out to be thin: a bbox-coordinate bug (commit \texttt{3d62e6a}), a double-scaling fix (\texttt{7eb7cc8}), and a deleted calibration file. The system came up because the ladder had caught everything else.

But the team-working element I need to name, because it is the one the rubric most requires, is \textbf{the wk17 FSM conflict and how I managed it.}

I had proposed an 11-state FSM for \texttt{main.py} with explicit transitions, timeout guards, and a geofence repulsor. Robin and Demetro pushed back in the same Teams thread. Robin argued that a two-out-of-five maintainability floor meant the FSM was over-engineered --- only two of five team members could read the code, and a team artefact that only the author can maintain is a single point of failure. Demetro argued the extra states would not survive a hardware-debug session with Edward, where you need to restart and reconfigure quickly. My first reaction was defensive. I genuinely believed the FSM was safer given our NFZ edge cases, and commit \texttt{ce5c036} later proved at least part of that belief correct when the RTL mode-tuple fix needed an explicit state transition to \texttt{LANDING}.

But I took 24 hours --- what Sch\"on (1983) calls reflection-in-action, though honestly it was closer to reflection-\textit{after}-action because I needed the overnight pause --- reread Robin's flight script, and realised half of my insistence was that the FSM was \textit{my} code. I was conflating ``is this the right design'' with ``is this the right design \textit{for this team},'' and these two kinds of rightness can come apart. I now see this is because Lencioni's Level 2, fear of conflict, produces exactly this pattern: when trust is low, technical disagreements get processed as identity threats, and the team defaults to ``artificial harmony'' rather than resolving the actual issue.

I proposed a compromise on the wk17 call: I would keep the full FSM for the autonomous path in \texttt{main.py}, but strip it to a lightweight polling structure exposed over HTTP for Robin's script. That is what the \texttt{--robin} flag in \texttt{passive\_watch\_2.py} actually is: not a feature, but an interface concession. Robin accepted, Demetro signed off, and the wk20 integration took an afternoon instead of a week.

The Hatton Level 4 move I want to make explicit: I was right about the FSM and wrong about the context. The technical design was sound --- the RTL bug later proved it. But imposing a sound design on a team that cannot maintain it is not engineering leadership; it is engineering narcissism. Graham (2025) puts this precisely in the innovation tools lecture: ``the tools are there to support, not replace, your thinking.'' My FSM was a tool; I had let it replace the team's thinking about what they could sustain.

\textbf{The second conflict: wk20 workload and communication.} A quieter but equally important conflict was the wk20 workload-and-comms tension with our PM after three flight days had been cancelled. Instead of confronting him directly --- something I historically do badly, because my instinct is to prepare a case rather than have a conversation --- I drafted a complaint letter in measured factual prose, iterated the tone six times across commits \texttt{26e3255}, \texttt{f3401e4}, \texttt{d946838}, \texttt{5af87d0}, \texttt{847518e}, and \texttt{4b0fd91}. I shared it only after Robin cut three sentences that were ``technically accurate but personally sharp.'' The PM disagreed with one paragraph; we rewrote it together; the letter went out with all five names.

The mediating action was \textit{externalising an internal conflict}: I turned a team tension into a written artefact that could be discussed, edited, and signed collectively, rather than letting it fester as unspoken resentment. I had not previously noticed that I can only handle confrontation when it is mediated by writing, and this is a pattern across my personal life, not just this project. Whether that is a cultural habit (Ukrainian directness channelled through bureaucratic form) or a personal one, I cannot fully separate. But I now recognise it as a tool with limitations: it works when there is time for iteration, and it fails when the issue needs a 30-second corridor conversation.

\textbf{The third conflict: team advocacy toward the course organiser.} The complaint letter itself (\texttt{docs/COMPLAINT\_LETTER.md}) was a team-level conflict with external authority. I want to name it because it demonstrates a different kind of conflict management: not resolving a disagreement within the team, but representing the team's collective grievance to an external stakeholder.

The letter documented, with factual precision and deliberate restraint, that ``three flight days were scheduled; our team completed zero successful flights\ldots approximately eight weeks\ldots zero seconds of real flight data.'' The tone was notably measured: ``We are not assigning personal blame. We are documenting what happened.'' I chose to lead this communication because I was the team member with the clearest view of what the hardware failures had cost in engineering terms, and because I had the written evidence (commit history, session logs) to back every claim.

The cross-team awareness I exercised: in the letter, I explicitly noted that borrowing the green team's drone ``does not constitute a plan'' because it risks ``disruption to a team whose hardware actually works.'' I was considering impact on other teams, not just my own. That sentence was the hardest to write, because the expedient answer was to grab the drone and go --- but the professional answer was to acknowledge the externality.

% ============================================================
% 2c. WHAT I WOULD CHANGE
% ============================================================

\subsection{2c. What I would change in a similar project}

Three concrete structural changes, each with a falsifiable signal.

\textbf{Change 1: Weekly decision log from wk1.} I would institute a one-line decision log from the first week --- every call that bound more than one person's work, with date, owner, and reason. The wk16 silent drift would have been caught by any written record of ``who owns what this week.'' The signal that it was working: zero overlapping-work surprises at integration checkpoints. Graham warns that teams which ``don't bother to think about how they work, possibly because they hate talking about management, will have to spend even more time trying to manage the team.'' That was us.

\textbf{Change 2: Front-loaded interface contracts.} I would front-load interface contracts \textit{before} parallel work begins --- specifically a two-hour whiteboard session in the first week of any new module to agree inputs, outputs, and failure modes at every seam. A single wk14 session on the \texttt{passive\_watch} HTTP contract would have removed the entire wk17 conflict. The signal: no integration surprises where two modules disagree on data format.

Drawing on the down-selection tools from Graham's (2025) innovation lecture, I would also propose using a simple MCDA matrix for architecture decisions. Our ROS vs. MAVLink decision, for example, was settled by informal discussion. A formal weighted comparison --- with criteria like Pi compatibility, learning curve, community support, and latency --- would have produced a more defensible decision and forced the team to articulate trade-offs explicitly rather than deferring to whoever argued loudest. Graham's warning resonates: ``if you skip the requirements analysis, you'll most likely find that the solution selection takes much longer anyway as the team will find it harder to agree on what good is supposed to look like.''

\textbf{Change 3: Documentation paired with conversation.} I would stop treating documentation as a substitute for conversation. The two pieces of writing that actually changed teammates' behaviour in this project were the 400-word pivot message (wk18 team call) and the 2000-word complaint letter --- both worked because they were short, timely, and accompanied by a live call. The 3000-word \texttt{VISION\_PIPELINE\_FOR\_COLLEAGUE.md}, by contrast, sat unread until Robin finally needed it three weeks later. The falsifiable signal: any document longer than 500 words must be paired with a 15-minute walkthrough within one week of being written, or I will treat it as unsent.

\ebi\ The deeper lesson, which connects to our Tuckman progression, is about how much time to invest in the team versus the technical work. Graham presents a ``crude model'' showing that good teams spend \textit{more} time managing themselves early and \textit{less} later, while poor teams do the opposite. We were the poor-team pattern: we spent almost no time on team management in wk12--16 (optimistic Forming), then spent weeks managing the fallout of overlapping work, miscommunication, and unclear ownership (extended Storming). If I could restart, I would advocate for one full afternoon in wk12 dedicated to nothing but team process: decision-log template, interface contracts, communication cadence, and a shared understanding of what ``done'' means for each module. That afternoon would have cost us five hours of technical work and saved us roughly thirty hours of rework.

Applying Lencioni's antidotes: I would start with a personal histories exercise in the first meeting --- simple sharing about backgrounds, working styles, and what each person needs to do their best work. The point is not to become friends; it is to build enough trust that wk17-style pushback feels normal, not threatening. And I would institute what Lencioni calls ``publication of goals and standards'' --- a shared document listing what each person has committed to deliver, by when, with what quality bar. The reason I didn't do this is that it felt bureaucratic and paternalistic. The reason I should have is that, as Lencioni puts it, ``peer pressure in the team, providing there's a foundation of trust, is way more effective than a bureaucratic measurement system.'' We had neither peer pressure nor a measurement system. We had me building things and hoping the team would notice.

The honest summary of 2c: I would change the team's investment in its own processes, and the person I would need to convince first is myself, because my instinct is always to write code rather than write a decision log, and every lesson in this section is a lesson about that instinct being wrong.
